% Generated draft: requires mathematical and source review.
\documentclass[11pt]{article}
\usepackage[margin=25mm]{geometry}
\usepackage{fontspec}
\setmainfont{Latin Modern Roman}
\usepackage{amsmath,amssymb,mathtools,unicode-math}
\setmathfont{Latin Modern Math}
\usepackage{enumitem,needspace,xurl,hyperref}
\setlength{\parindent}{0pt}
\setlength{\parskip}{5pt}
\setlength{\emergencystretch}{3em}
\newcommand{\N}{\mathbb{N}}
\newcommand{\Z}{\mathbb{Z}}
\newcommand{\Q}{\mathbb{Q}}
\newcommand{\R}{\mathbb{R}}
\newcommand{\C}{\mathbb{C}}
\newcommand{\F}{\mathbb{F}}
\newcommand{\A}{\mathbb{A}}
\newcommand{\PP}{\mathbb{P}}
\newcommand{\E}{\mathbb{E}}
\newcommand{\eps}{\varepsilon}
\newcommand{\dd}{\,\mathrm{d}}
\begin{document}
\section*{UnsolvedMath ID: 800129}
\textbf{Source problem code:} catalogue entry
\subsection{Definitions and mathematical statement}
Let \(p\) be a prime, and let \(\mathbb F_p\) denote the field with \(p\) elements. Let \(M\) be a PL manifold, subject to the unresolved scope qualifications below. Write \(\mathrm{PL}(M)\) for the group of piecewise-linear self-homeomorphisms of \(M\), equipped with the topology intended in the source, and write \(\mathrm{PL}(M)^\delta\) for the same group with the discrete topology. For a topological group \(H\), let \(BH\) denote its classifying space. The identity homomorphism \(\mathrm{PL}(M)^\delta\to\mathrm{PL}(M)\) induces a comparison map
\[
\iota_M\colon B\mathrm{PL}(M)^\delta\longrightarrow B\mathrm{PL}(M).
\]
The question is whether, for every integer \(k\geq 0\), the induced map on singular homology
\[
(\iota_M)_*\colon H_k\bigl(B\mathrm{PL}(M)^\delta;\mathbb F_p\bigr)\longrightarrow H_k\bigl(B\mathrm{PL}(M);\mathbb F_p\bigr)
\]
is an isomorphism.

The supplied statement does not specify the topology on \(\mathrm{PL}(M)\), the dimension or compactness assumptions on \(M\), whether boundary is allowed, or any boundary-fixing requirement. It also does not explicitly state the range of manifolds or primes over which the question is quantified. Thus the displayed question retains these unresolved parameters rather than specifying a fully determined universal conjecture.

The source presents this as an analogue of Milnor's conjecture: for a finite-dimensional Lie group \(G\), the comparison map
\[
H_k(BG^\delta;\mathbb F_p)\longrightarrow H_k(BG;\mathbb F_p)
\]
is conjectured to be an isomorphism for all \(k\geq 0\). As motivation, the source tentatively suggests that Greenberg's work on PL foliations indicates an inductive construction of the classifying space of PL Haefliger structures from classifying spaces of discrete Lie groups. This motivation is recorded as a source qualification, not asserted here as a precise theorem.
\subsection{Short English statement}
The question asks whether making the group of PL self-homeomorphisms of \(M\) discrete preserves the homology of its classifying space with coefficients in \(\mathbb F_p\), in every degree. The intended topology, manifold hypotheses, boundary conventions, and scope of quantification remain unspecified in the supplied record.
\subsection{Sources}
\begin{itemize}
\item Local UnsolvedMath registry, permanent ID \(800129\), source problem code \texttt{catalogue entry}, title: \emph{Foliations and Diffeomorphism Groups --- Question 70}. The supplied record attributes the problem to Nariman and Yazdi, with proposed year 2025, and labels its status as open; this status is supplied metadata, not an independently verified assessment.
\item Nariman and Yazdi, \emph{Problem List on Foliations and Diffeomorphism Groups} (2025), Question 70, \url{https://arxiv.org/abs/2503.18274}. The supplied provenance records access on 2026-10-16 and extraction by \texttt{arxiv\_latex}; these provenance details are not independently verified here.
\item The supplied statement attributes the finite-dimensional Lie-group comparison conjecture to Milnor and mentions Greenberg's work on PL foliations as tentative motivation. It supplies no separate bibliographic references for these attributions.
\end{itemize}
\end{document}
